\documentclass[twoside]{report}
\usepackage{german,epsf,haupt}
\dqwarninglevel{0}
\usepackage{amsfonts}
\newcommand{\C}{\mathbb{C}}
\newcommand{\R}{\mathbb{R}}

\newcommand{\includeeps}[2]{
   \def\epsfsize##1##2{#2##1}
   \centerline{\epsffile{#1}}}

\begin{document}
\begin{titlepage}\Large
\begin{center}
{\LARGE\bf Seminar} \\[1cm]
{\huge\bf Quantum Computing\par} 
\vspace*{1cm}
\unilogo{50}\\[1cm]
{\bf Universit"at Karlsruhe (TH)}\\
{Fakult"at f"ur Informatik}\\
{\em Institut f"ur Algorithmen und Kognitive Systeme}
\vfill
Prof.~Dr.~Th.~Beth\\
Dipl.~Inform.~M.~Grassl\\
Dipl.~Inform.~J.~M"uller--Quade
\vfill\vfill 
Sommersemester 1996
\vfill
\vfill
\end{center}
\end{titlepage} 
 
\newpage
\thispagestyle{empty}
\ 
\newpage 

\pagenumbering{roman}
\tableofcontents
\cleardoublepage
\pagenumbering{arabic}
%%
%% hier wird die eigentlich Ausarbeitung eingebunden (eigenes File,
%% mit \input
%% \documentstyle[german,11pt,a4,mathsym]{article}
%% \begin{document}
\newcommand{\scalar}[3]{\left< #1 | #2 | #3 \right>}
\newcommand{\betrag}[1]{\left| #1 \right|}
\newcommand{\norm}[1]{\| #1\|}
\newcommand{\raum}{H}
%\newcommand{\traum}[1]{$H_{#1}$}
\newcommand{\quadr}[1]{\left( #1 \right) ^2}
\newcommand{\menge}[1]{\left\{ #1 \right\}}
\newcommand{\scalars}[2]{\left< #1 | #2 \right>}
\newcommand{\tscalars}[6]{ \left< #1 #3 | #2 #4 \right>_{#5 , #6} }
\newcommand{\Ima}{\mbox{Im}}
\newcommand{\Spur}[1]{\mbox{Sp} \left( #1\right)}
\newcommand{\Proj}[1]{\ket{#1}\bra{#1}}
\newcommand{\Prob}[2]{{\cal P}_{#1} \left( #2 \right)}
\newcommand{\bra}[1]{\left< #1 \right|}
\newcommand{\ket}[1]{\left| #1 \right>}
\newcommand{\tbra}[2]{\left< #1 \right|_{#2}}
\newcommand{\tket}[2]{\left| #1 \right>_{#2}}
\newcommand{\mean}[1]{\left< #1 \right>}
\newcommand{\var}[1]{\left( \triangle \left< #1 \right> \right)^2}
\newcommand{\unscharf}[1]{ \triangle \left< #1 \right>}
\newcommand{\id}{\hbox{\small1\kern-3.8pt\normalsize1}}
\newcommand{\difft}{{\partial \over \partial t}}
\newcommand{\kommutator}[2]{\left[ #1 , #2 \right]}
\newcommand{\Hamilton}{{\cal H}}
\newcommand{\Poisson}[2]{\lbrace #1 , #2 \rbrace}
\newcommand{\tupel}[1]{\left( #1 \right)}
\Vortrag{Quantentheorie}{Pawel Wocjan}

{\bf Vorbemerkung} \\
Im ersten Teil wird eine Einf"uhrung in die Grundgedanken der Quantentheorie gegeben und ihre mathematische Formulierung erarbeitet. Im zweiten Teil wird der
bisher entwickelte Kalk"ul um statistische Methoden erweitert. Diese Erweiterung wird notwendig, wenn nur eine unvollst"andige Information "uber den Zustand des
Quantensystems vorliegt. 

\section{Formulierung und Interpretation der Quantentheorie}
  \subsection{Der Zustandsraum}
    Die mathematische Theorie der Hilbertr"aume stellt das richtige Werkzeug dar, um Quantensysteme zu beschreiben. Jedes Quantensystem wird in einem geeigneten
    komplexen Hilbertraum, dem Zustandsraum \raum, beschrieben. Dieser Zustandsraum \raum{} wird durch die Struktur des speziell betrachteten physikalischen Problems
    bestimmt. Im folgenden wird die Korrespondenz zwischen den mathematischen Objekten in der Theorie und den physikalischen Gegebenheiten eines Quantensystems 
    beschrieben.
    
  \subsection{Zust"ande} 
      Die Einheitsvektoren $\ket{\Psi} \in$ \raum{} hei\3en Zustandsvektoren. Zwei Zustandsvektoren $\ket{\Psi}$ und $\ket{\Phi}$ hei\3en "aquivalent, wenn 
      $\ket{\Psi} = e^{i\Theta}\ket{\Phi}$ f"ur $\Theta \in \R$ gilt. 
      Jeder physikalische Zustand eines Quantensystems entspricht einem Zustandsvektor $\ket{\Psi}$ in \raum. Dabei repr"asentieren "aquivalente Zust"ande denselben
      physikalischen Zustand des Systems. \\
      Betrachten wir als Beispiel die lineare Ionenfalle, um die Unterschiede zwischen der klassischen und der quantentheoretischen Beschreibung zu sehen. Die lineare Ionenfalle
      ist eine der bisher
      erwogenen Hardwarel"osungen f"ur einen Quantencomputer. Sie ist (in einfachster Beschreibungsstufe) ein System mit $n$ Komponenten, die jeweils zwei Zust"ande $\ket{0}$ und
      $\ket{1}$ annehmen k"onnen. Klassisch wird der Zustand dieses Systems mit $n$ Bits vollst"andig beschrieben; jeder m"ogliche Zustand des Systems entspricht einer Ecke eines
      $2^n$-dimensionalen Hyperw"urfels. Quantentheoretisch dagegen sind f"ur eine vollst"andige Beschreibung $2^n-1$ komplexe Zahlen notwendig. Der Zustand des Quantensystems
      entspricht einem Zustandsvektor in einem $2^n$ dimensionalen Hilbertraum, dem Zustandsraum \raum. Jeder Zustand dieses Systems zu einem beliebigen Zeitpunkt wird durch einen
      Einheitsvektor
      beschrieben. Die Zust"ande liegen auf also einer $2^n$-dimensionalen Einheitskugel. Wenn man einen Zustandsvektor mit einem beliebigen
      Phasenfaktor $e^{i\Theta}, \Theta \in \R,$ multipliziert, erh"alt man einen "aquivalenten Zustandsvektor. Man braucht also nur $2^n-1$ komplexe Zahlen (Charakterisierung
      der L"ange und Phase der Basisvektoren relativ zueinander), um den Zustand eindeutig festzulegen, da ein globaler Phasenfaktor physikalisch unrelevant ist. Man kann diesen
      an sich beliebigen Phasenfaktor eliminieren, wenn man die Zust"ande des Quantenssystems nicht mit Zustandsvektoren $\ket{\Psi}$, sondern mit den Projektionsoperatoren
      $\Proj{\Psi}$ beschreibt, weil sich durch Adjugieren der Phasenfaktor aufhebt
      $$ e^{i\Theta}\ket{\Psi} \mbox{"aquivalent zu} \ket{\Psi} \Longrightarrow \ket{\Psi}e^{-i\Theta}e^{i\Theta}\bra{\Psi} = \Proj{\Psi}. $$
      Im zweiten Teil wird genau dieses Konzept benutzt, um die reinen Zust"ande und statistische Mischungszust"ande einheitlich mit Hilfe solcher Operatoren, den Dichteoperatoren,
      zu beschreiben. Ein beliebiger Zustand des Systems kann als eine lineare Superposition der Basiszust"ande $\ket{i} = \ket{\mbox{Bin"ardarstellung von $i$}}$
      $$ \ket{\Phi} = \sum_{i=0}^{2^n-1} a_i \ket{i}$$
      dargestellt werden, wobei f"ur die Koeffizientenn $a_i \in \C$ gilt $\sum_{i} \betrag{a_i}^2 = 1$. Zu jedem klassisch m"oglichen Zustand gibt es
      einen Basisvektor des Zustandsraums \raum. 
   \subsection{Observablen}
      Bei der Messung an Quantensystemen ist im Gegensatz zu den makroskopischen Systemen, die korrekt durch die klassische Physik beschrieben werden, eine
      Abstraktion von der Me\3vorrichtung prinzipiell nicht mehr m"oglich. Man ist nicht
      imstande, irgendwelche Observable des Quantensystems gleichzeitig zu messen. Nur der Me\3wert derjenigen Observable, die man gerade gemessen hat, kann
      dem Quantensystem zugeordent werden. Die Me\3werte anderer Observablen d"urfen also im allgemeinen nicht aus den Me\3werten anderer Observabler berechenbar 
      sein. Es mu\3 daher ein mathematisches Kalk"ul entwickelt werden, der diese Tatsache ber"ucksichtigt. Das bedeutet, da\3 quantentheoretisch die Observablen in
      anderer Weise zu erfassen sind als in der klassischen Physik. Die Operatoren des Zustandsraumes \raum{} des Quantensystems, die bestimmte Eigenschaften erf"ullen, sind die 
      geeigneten mathematischen Objekte f"ur die quantentheoretischen Beschreibung der Observablen.
      Die im allgemeinen gegebene Nichtvertauschbarkeit zweier Operatoren erlaubt n"amlich automatisch die mathematische Erfassung der nicht beliebigen scharfen, gleichzeitigen 
      Me\3barkeit zweier Observablen. Da es sich ferner erweisen wird, da\3 die Eigenwerte der Operatoren die m"oglichen Me\3werte der Observablen sind, beschr"ankt
      man sich auf hermitesche Operatoren, weil deren Eigenwerte reell sind, wie man es von physikalischen Me\3werten  verlangt. Au\3erdem m"ussen die Eigenvektoren des 
      Operators eine Basis des Zustandraumes \raum{} bilden. 
   \subsection{Messungen}
      Um "uber den Ausgang der Messung einer quantenmechanischen Observablen zu einem bestimmten Zeitpunkt Aussagen machen zu k"onnen, ist es notwendig, den
      Zustand, in dem sich das physikalische System unmittelbar vor der Messung befindet, durch einen Zustandsvektor zu beschreiben. Durch eine mathematische
      Verkn"upfung zwischen dem Zustandsvektor und dem Operator, der die zu messende Observable darstellt, mu\3 es dann m"oglich sein, die gew"unschten 
      quantentheoretischen Aussagen "uber den Me\3vorgang zu machen. 
      \subsubsection{Erwartungswert}
         Im Gegensatz zur klassischen Physik ergeben sich in der Quantentheorie bei mehrmaliger Ausf"uhrung einer Messung derselben Observablen im allgemeinen
         verschiedene Me\3ergebnisse, auch wenn das System vor der Messung immer im gleichen Zustand war. Wenn man also wei\3, da\3 sich das System im Zustand
         ${\ket{\Phi}}$ befindet, so kann man "uber den Ausgang der Messung einer Observablen $A$ nur Vorhersagen "uber den Mittelwert oder den Erwartungswert
         $\mean{A}$ der verschiedenen m"oglichen Me\3werte machen. Der radikale Unterschied zwischen der klassischen Physik und der Quantentheorie ist derjenige, da\3 der 
         Me\3vorgang an Quantensystemen nicht deterministisch ist. Dies ist eine inherente Eigenschaft der Natur, die prinzipiell nicht umgangen werden kann. Die fundamentale
         Behauptung der Quantentheorie (Born) ist, da\3 dieser Erwarungswert durch das
         Skalarprodukt $\scalar{\Phi}{A}{\Phi}$ gegeben ist. Dieser Zusammenhang zwischen dem Me\3vorgang und den im mathematische Kalk"ul der Quantentheorie 
         vorkommenden Objekten best"atigt sich in allen Experimenten. \\
         Die statistische Aussagen der Quantentheorie lautet also:
         Der Zustand eines Quantensystems wird durch einen Zustandsvektor $\ket{\Phi}$ im geeigneten Zustandsraum \raum{} beschrieben. F"uhrt man an dem System die Messung einer
         Observablen $A$ sehr oft aus, und ist vor jeder Messung das System immer im gleichen Zustand $\ket{\Phi}$, so ist der Erwartungswert $\mean{A}$ der 
         Me\3ergebnisse gegeben durch
         $$ \mean{A} = \scalar{\Phi}{A}{\Phi}. $$
         Sind die Eigenwerte $\alpha_i$ von $A$ alle einfach, so gilt mit Hilfe der Entwicklung von $\ket{\Phi}$ nach den Eigenvektoren $\menge{\ket{\Psi_i}}$ von $A$
         $$ \mean{A} = \sum_i \alpha_i\betrag{\scalars{\Psi_i}{\Phi}}^2. $$
         Sind nicht alle Eigenwerte $\alpha_i$ von $A$ einfach, so erh"alt man mit der Enwicklung nach den Eigenvektoren $\menge{\ket{\Psi_{ij}}}$
         $$ \mean{A} = \sum_i \alpha_i \sum_j \betrag{\scalars{\Psi_{ij}}{\Phi}}^2, $$
         wobei $\ket{\Psi_{ij}}$ ein Eigenvektor zum Eigenwert $\alpha_i$ ist.
      \subsubsection{Varianz}
         F"uhrt man die Messung einer Observablen $A$ mehrmals aus (wobei das System vor der Messung immer im gleichen Zustand sei), so streuen im allgemeinen die
         Me\3werte um den Mittelwert $\mean{A}$. Das Ma\3 f"ur die Streuung ist ihre Varianz $\var{A}$, die als Erwartungswert der quadratischen Abweichung von $\mean{A}$
         definiert ist.
         \begin{eqnarray*}
            \var{A} & = & \mean{\left( A-\mean{A}\id\right)^2} \\
                    & = & \scalar{\Phi}{\left( A-\mean{A}\id\right)^2}{\Phi}
         \end{eqnarray*}
         Wegen
         $$ \mean{\left( A-\mean{A}\id\right)^2} = \mean{A^2} - 2\mean{A\mean{A}} + \mean{\mean{A}^2} = \mean{A^2} - \mean{A}^2 $$
         gilt auch
         $$ \var{A} = \mean{A^2} - \mean{A}^2 = \scalar{\Phi}{A^2}{\Phi} - \scalar{\Phi}{A}{\Phi}^2 $$
         Mit $A$ ist auch $A-\mean{A}\id$ hermitesch, so da\3 man f"ur die Varianz auch
         $$ \var{A} = \norm{\left(A-\mean{A}\id\right)\ket{\Phi}}^2 \ge 0 $$
         schreiben kann. \\
         Die Unsch"arfe $\unscharf{A} = \sqrt{\var{A}}$ ist nichts anderes als die L"ange des Vektors $\left(A-\mean{A}\id\right)\ket{\Phi}$,
         $$ \unscharf{A} = \norm{\left(A-\mean{A}\id\right)\ket{\Phi}}. $$
       \subsubsection{Der Einflu\3{} der Messung}
         Es stellt sich die Frage, bei welchen speziellen Zustandsvektoren $\ket{\Phi_A}$ die Me\3werte bei der Messung von $A$ streuungsfrei sind. Dazu ist notwendig und
         hinreichend, da\3 \\ $\norm{\left(A-\mean{A}\id\right)\ket{\Phi_A}} = 0$ oder "aquivalent dazu
         $$ A\ket{\Phi_A} = \mean{A}\ket{\Phi_A} $$
         gilt. Der Zustandsvektor $\ket{\Phi_A}$ mu\3 also ein Eigenvektor von $A$ und $\mean{A}$ der zugeh"orige Eigenwert sein. Da bei Streuugsfreiheit der 
         Erwartungswert zu einem mit Sicherheit vorliegenden Me\3wert wird, kann man schlie\3en: \\
         Ergibt sich bei einer Messung von $A$ ein Me\3wert mit Sicherheit, so ist dieser Me\3wert ein Eigenwert von $A$, und der Zustandsvektor vor der Messung mu\3
         der zugeh"orige Eigenvektor gewesen sein. Nat"urlich gilt auch die Umkehrung dieses Satzes: \\
         Ist der Zustandsvektor $\ket{\Phi}$ ein Eigenvektor, so ergibt sich bei der Messung von $A$ mit Sicherheit der zugeh"orige Eigenwert als Me\3wert. Ist der
         Eigenwert entartet, so gilt Entsprechendes, wenn der Zustandsvektor im zugeh"origen Eigenraum liegt. \\
         Es mu\3 noch gekl"art werden, welcher Zustandsvektor unmittelbar nach der Messung einer Observablen vorliegt. Wenn die Messungen "uberhaupt einen 
         physikalischen Sinn haben sollen, so ist zu fordern, da\3 bei einer unmittelbaren Wiederholung derselben Messung sich das Ergebnis der vorangehenden Messung nicht
         "andert. Bei der zweiten Messung mu\3 sich also mit Sicherheit dasselbe Ergebnis wie in der ersten ergeben. Deswegen m"ussen  also auch bei einem beliebigen                Anfangszustand $\ket{\Phi}$ vor der Messung die "uberhaupt m"oglichen Me\3werte einer Observablen $A$ ihre reellen Eigenwerte sein. Der Zustandsvektor geht
         aufgrund der Einwirkung der Messung in den zugeh"origen Eigenvektor von $A$ "uber, der zu dem gemessenen nichtentarteten Eigenwert geh"ort. Dies wird
         als die Reduktion des Zustandsvektors bezeichnet. \\
         Wenn bei der Messung einer Observablen $A$ ein entarteter Eigenwert $\alpha_i$ eintritt, so mu\3 der Zustandsvektor $\ket{\Phi_{\alpha_i}}$ nach der Messung im
         Eigenraum $E_{\alpha_i}$ liegen. Wird dieser von den orthonormierten Eigenvekoren $\menge{\ket{\Psi_{ij}}}$, wobei $j$ der Index des Eigenvektors im Eigenraum angibt,
         aufgespannt, so gilt also
         $$ \ket{\Phi_{\alpha_i}} = \sum_j \ket{\Psi_{ij}} c_{ij}. $$
         Es kann gezeigt werden, da\3  die Koeffizienten $c_{ij}$ nicht wilk"urlich, sondern proportional zu $\scalars{\Psi_{ij}}{\Phi}$ sind
         $$ \ket{\Phi_{\alpha_i}} = c \sum_j \ket{\Psi_{ij}}\scalars{\Psi_{ij}}{\Phi}.$$
         Der Zustandsvektors $\ket{\Phi_{\alpha_i}}$ entsteht aus $\ket{\Phi}$ durch Projektion auf den Eigenraum $E_{\alpha_i}$. Bestimmt man den Betrag der 
         Konstanten $c$ aus der Normierungsbedingung $\scalars{\Phi_{\alpha_i}}{\Phi_{\alpha_i}} = 1$, so ergibt sich
         $$ \ket{\Phi} \Longrightarrow \frac{\sum_j \ket{\Psi_{ij}} \scalars{\Psi_{ij}}{\Phi}} {\sqrt{\sum_j\scalars{\Psi_{ij}}{\Phi}}} $$
         Bei dem Me\3vorgang werden also gerade die Komponeten von $\ket{\Phi}$ herausgefiltert, die auf dem Eigenraum $E_{\alpha_i}$ senkrecht stehen.
       \subsubsection{Die Wahrscheinlichkeit f"ur das Eintreten eines Me\3wertes}
         Vor der Ausf"uhrung der Messung einer Observablen $A$ l"a\3t sich nicht mit Sicherheit sagen, welcher Me\3wert sich aus der Vielfalt der Eigenwerte ergibt -
         abgesehen von dem Fall, da\3 vor der Messung der Zustandsvektor $\ket{\Phi}$ ein Eigenvektor von $A$ ist. Es sind vielmehr bei dem beliebigen Zustandsvektors
         $\ket{\Phi}$ nur Wahrscheinlichkeitsaussagen m"oglich. Durch Vergleich der Definition des Erwartungswertes
         $$ \mean{A} = \sum_i \Prob{\ket{\Phi}}{\alpha_i} \alpha_i $$
         mit der quantentheoretischen
         $$ \mean{A} = \sum_i \betrag{\scalars{\Psi_i}{\Phi}}^2 \alpha_i $$
         ergibt sich f"ur die Wahrscheinlichkeit, da\3 der Eigenwert $\alpha_i$ gemessen wird,
         $$\Prob{\ket{\Phi}}{\alpha_i} = \betrag{\scalars{\Psi_i}{\Phi}}^2 .$$
         Man nennt daher die Komponenten $\scalars{\Psi_i}{\Phi}$ des Zustandsvektors $\ket{\Phi}$ Wahrscheinlichkeitsamplituden f"ur den Me\3wert $\alpha_i$.
         Die Wahrscheinlichkeit $\Prob{\ket{\Phi}}{\alpha_i}$ hat die Gestalt $\scalars{\Phi}{\Psi_i}\scalars{\Psi_i}{\Phi}$ und kann daher als Erwartungswert des
         Projektionsoperators $P_{\ket{\Psi_i}} = \Proj{\Psi_i}$ aufgefa\3t werden,
         $$ \Prob{\ket{\Phi}}{\alpha_i} = \mean{P_{\ket{\Psi_i}}} = \scalar{\Phi}{P_{\ket{\Psi_i}}}{\Phi}.$$
         Dem Projektionsoperator $P_{\ket{\Psi_i}}$ entspricht die Observable "{}Ist das System im Zustand $\ket{\Psi_i}$?". Da die Eigenwerte von Projektionsoperatoren nur
         1 oder 0 sind, bedeuten diese Me\3werte, da\3 die Messung entscheidet, ob diese Frage mit Ja oder Nein zu beantworten ist. Vor der Messung kann man aber im allgemeinen 
         nicht mit Sicherheit vorhersagen, welche Antwort durch die Messung gegeben wird. Man kann vielmehr nur den Erwartungswert $\mean{P_{\ket{\Psi_i}}}$ angeben, der sich
         einstellt, wenn man - immer wieder vom gleichen Zustand $\ket{\Phi}$ ausgehend - die Frage h"aufig durch die Messung beantworten l"a\3t. Der Erwartungswert von
         $P_{\ket{\Psi_i}}$ ist aber gerade der Mittelwert "uber viele Entscheidungen "{}Ist das System in $\ket{\Psi_i}$?" gedanklich gleich der Wahrscheinlichkeit 
         $\Prob{\Phi}{\alpha_i}$ f"ur das Eintreffen von $\alpha_i$. \\
         Bei entarteten Eigenwerten gewinnt man die Wahrscheinlichkeit $\Prob{\Phi}{\alpha_i}$ f"ur das Eintreffen eines Eigenwerts $\alpha_i$ aus der Gleichung
         $$ \mean{A} = \sum_i \alpha_i \sum_j \betrag{\scalars{\Psi_{ij}}{\Phi}}^2 $$
         f"ur den Erwarungswert. 
         \begin{eqnarray*}
            \Prob{\Phi}{\alpha_i} & = & \sum_j \betrag{\scalars{\Psi_{ij}}{\Phi}}^2 \\
                                  & = & \sum_j \scalars{\Phi}{\Psi_{ij}}\scalars{\Psi_{ij}}{\Phi} \\
                                  & = & \sum_j \mean{P_{\ket{\Psi_{ij}}}}
         \end{eqnarray*}
         Diese Wahrscheinlichkeit kann mit der Spur als
         \begin{eqnarray*}
           \Prob{\Phi}{\alpha_i} & = & \sum_j \scalars{\Phi}{\Psi_{ij}}\scalars{\Psi_{ij}}{\Phi} \\
                                 & = & \Bigl< \Phi \Big| \bigl( \sum_j \Proj{\Psi_{ij}} \bigr) \Big| \Phi \Bigr> \\
                                 & = & \scalar{\Phi}{P_{\alpha_i}}{\Phi} \\
                                 & = & \Spur{ P_{\ket{\Phi}} P_{\alpha_i} }
         \end{eqnarray*}
         ausgedr"uckt werden. Sie kann als wieder als Erwartungswert des Projektionsopetartors $P_{\alpha_i}$ auf den Eigenraum $E_{\alpha_i}$ aufgefa\3t werden.
         
       \subsubsection{Die Unsch"arferelation}
         Kommutieren zwei Observable A und B nicht (ihr Kommutator $\kommutator{A}{B} = AB - BA \neq 0$), so existiert im allgemeinen kein Zustand, f"ur den die Unsch"arfe von
         $A$ und $B$ gleichzeitig Null ist. Es mu\3 als eine Beziehung zwischen den Unsch"arfen $\unscharf{A}$ und $\unscharf{B}$ bestehen, die das gleichzeitige
         Verschwinden von $\unscharf{A}$ und $\unscharf{B}$ ausschlie\3t, die Unsch"arferelation. Es gilt:
         $$ \unscharf{A}=\norm{\left( A-\mean{A}\id\right)\ket{\Phi}} = \sqrt{\scalar{\Psi}{\left( A-\mean{A}\id\right)^2}{\Psi}} $$
         $$ \unscharf{B}=\norm{\left( B-\mean{B}\id\right)\ket{\Phi}} = \sqrt{\scalar{\Psi}{\left( B-\mean{B}\id\right)^2}{\Psi}} $$
         Mit Hilfe der Schwarzschen Ungleichung erh"alt man
         \begin{eqnarray*}
           \unscharf{A}\unscharf{B} & \ge & \betrag{\scalar{\Psi}{\left( A-\mean{A}\id\right)\left( B-\mean{B}\id\right)}{\Psi}} \\
                                    & \ge & \betrag{\Ima\scalar{\Psi}{\left( A-\mean{A}\id\right)\left( B-\mean{B}\id\right)}{\Psi}} \\
                                    & = &   \betrag{ \frac{1}{2i} \left(\scalar{\Psi}{\left( A-\mean{A}\id\right)\left( B-\mean{B}\id\right)}{\Psi} -
                                                                        \scalar{\Psi}{\left( B-\mean{B}\id\right)\left( A-\mean{A}\id\right)}{\Psi} \right) } \\
                                    & = &   \frac{1}{2}\betrag{\scalar{\Psi}{AB-BA}{\Psi}} \\
                                    & = &   \frac{1}{2}\betrag{\scalar{\Psi}{\kommutator{A}{B}}{\Psi}} \\
                                    & = &   \frac{1}{2}\betrag{\mean{\kommutator{A}{B}}_{\ket{\Psi}}}
         \end{eqnarray*}
         Diese Absch"atzung gilt ganz allgemein f"ur alle Observablen. Man erh"alt die Heisenbergsche Unsch"arferelation, indem man den Impulsoperator und den Ortsoperator in die 
         obige Formel einsetzt.

     \subsection{Ein vollst"andiger Satz kommutierender Observablen}
         Die Nichtkommutativit"at zweier Operatoren bringt die nicht gleichzeitige, scharfe Me\3barkeit der zugeh"origen physikalischen Observablen zum Ausdruck. Der
         Zustand eines Quantensystems ist damit nicht wie in der klassischen Physik durch die genaue Angabe der Me\3werte aller Observablen festzulegen. \\
         Ist der Eigenwert $\alpha$ der Observablen $A$ einfach, so wird durch die Messung dieser Observablen bei Einstellung des Me\3wertes $\alpha$ der 
         Zustand nach der Messung eindeutig (bis auf einen physikalisch uninteressanten Phasenfaktor) festgelegt: $\ket{\Psi'} = \ket{\Psi_{\alpha}}$. Besitzt ein 
         Operator $A$ nur einfache 
         Eigenwerte, so nennt man die zugeh"orige Observable vollst"andig, weil durch ihre Messung der Zustand des Quantensysstems vollst"andig festgelegt wird - 
         eine dar"uber hinausgehende, genauere Information "uber den Zustand des Systems erlaubt die Quantentheorie nicht. \\
         Sind die Eigenwerte des Operators $A$ nicht alle einfach, so wird durch die Messung eines seiner entarteten Eigenwerte $\alpha$ der Zustandsvektor
         $\ket{\Psi}$ auf den zugeh"origen Eigenraum $E_{\alpha}$ projiziert. Die Richtung des entstehenden Zustandsvektors $\ket{\Psi'}$ im Eigenraum h"angt also
         von der Richtung von $\ket{\Psi}$ vor der Messung ab. Kennt man diese Richtung nicht, so wei\3 man nach der Messung von $\alpha$ nichts "uber die Richtung
         von $\ket{\Psi'}$. Zur genaueren Festlegung von $\ket{\Psi'}$ ist es dann n"otig, neben der Observablen $A$ noch
         eine weitere, mit $A$ kommutierende Observable $B$ zu messen. Sind die simultanen Eigenr"aume von $A$ und $B$ alle eindimensional, so legt die gleichzeitige
         Messung von $A$ und $B$ den Zustandsvektor eindeutig fest: $\ket{\Psi''} = \ket{\Psi_{\alpha\beta}}$. Man nennt die beiden vertr"aglichen Observablen dann einen
         vollst"andigen Satz von Observablen; ihr simultane Messung hat dieselbe Wirkung wie die Messung eines vollst"andigen Operators. \\
         Sind die simultanten Eigenr"aume von $A$ und $B$ jedoch noch nicht eindimensional, so sind noch weitere, kommutierende Operatoren $C,\ldots$ hinzuzuf"ugen,
         bis schlie\3lich die simultanten Eigenr"aume all dieser Operatoren eindimensional sind. Durch die gleichzeitige Messung eines solchen vollst"andigen Satzes
         kommutierener Observablen wird der Zustand des Systems vollst"andig bestimmt,
         $$ \ket{\Psi'''^{\dots}} = \ket{\Psi_{\alpha\beta\gamma\ldots}}. $$ 
         
     \subsection{Zeitliche Entwicklung}
        Die Gleichung
        $$ \ket{\Psi\left( t\right)} = \exp\left( -\frac{itH}{\hbar}\right) \ket{\Psi_0} $$
        beschreibt die zeitliche Evolution des Quantensystems, wenn $\ket{\Psi_0}$ dem Zustand des Systems zum Zeitpunkt $t=0$ und $\ket{\Psi\left( t\right)}$
        zum Zeitpunkt $t$ entspricht und $H$ der Hamiltonoperator ist. 
        $$ \left\{ \exp\left( -\frac{itH}{\hbar}\right)\right\} $$
        ist die einparametrige unit"are Gruppe, die durch den schiefadjungierten Operator $ -\frac{itH}{\hbar}$ erzeugt wird. Den Evolutionsoperator
        $$ U\left( t\right) = \exp\left( -\frac{itH}{\hbar}\right) $$
        kann man aus der Schr"odingergleichung 
        $$ i\hbar {\partial \over \partial t} \ket{\Psi\left( t\right)} = H\ket{\Psi\left( t\right)} $$
        ableiten. Die Schr"odingergleichung ist eine Differentialgleichung erster Ordung in der Zeit, und deswegen ist $\ket{\Psi\left( t\right)}$ eindeutig durch den
        Anfangszustand $\ket{\Psi_0}$ bestimmt. Au\3erdem gen"ugt sie dem Superpositionsprinzip, da\3 hei\3t Linearkombinationen von L"osungen stellen wieder
        L"osungen dar. Deshalb k"onnen Interferenzeffekte auftreten. \\
        Da der Evolutionsoperator $U\left( t\right)$ f"ur alle $t$ unit"ar ist, gilt
        $$ \scalars{\Psi\left( t\right)}{\Psi\left( t\right)} = \scalars{\Psi_0}{\Psi_0} = 1 $$
        Bei der zeitlichen Evolution eines Quantensystems werden also Zust"ande in Zust"ande "uberf"uhrt. Dabei beschreibt der Zustandsvektor einen Weg auf der Einheitskugel. 
    
 
\section{Quantentheorie bei unvollst"andiger Information "uber den Zustand des Systems}
  \subsection{Analogie zur klassischen statistischen Mechanik}
    Beim Hamilton-Formalismus wird das System in der klassischen Mechanik durch die Angabe der 
    Orts- und Impulskoordinaten $q_i$, $p_i$ beschrieben, die einen Punkt im Phasenraum darstellen.
    Bei vielen Problemen sind jedoch diese Werte keineswegs exakt bekannt. Man kann in solchen
    F"allen lediglich sagen, da\3 sich das System in bestimmten Bereichen des Phasenraumes
    nicht (oder kaum) befindet, w"ahrend es in anderen Bereichen mit gro\3er Wahrscheinlichkeit
    vorkommt. Man dr"uckt also die Information, die man "uber das System (zu einem
    bestimmten Zeitpunkt) hat, durch eine Wahrscheinlichkeitsverteilung
    $$ \rho = \rho\left( q_1 ,\ldots ,p_1 \ldots \right) $$
    im Phasenraum aus. Sie ist die Wahrscheinlichkeitsdichte daf"ur, da\3 sich das
    System in einem Phasenpunkt befindet (klassische statistische Mechanik). Man
    normiert $\rho$, indem man das Integral "uber den ganzen Phasenraum eins setzt,
    $$ \int\rho\mbox{d}\Omega = 1 \mbox{, d}\Omega = \mbox{d}q_1\ldots \mbox{d}p_1\ldots $$
    In der Quantentheorie mu\3  eine vollst"andige Observable an dem System gemessen werden, um den Zustandsvektor im Zustandsraum \raum{} eindeutig zu pr"aparieren. Dies kann 
    durch die Messung eines vollst"andigen Satzes kommutierender Observablen geschehen. F"ur sp"atere Zeitpunkte kann der Zustandsvektor 
    mit Hilfe der Schr"odingergleichung berechnet werden. In vielen F"allen ist eine solche vollst"andige Messung "uberhaupt nicht zu realisieren. Es
    k"onnen lediglich hochgradig entartete Observable gemessen werden. Nach der Messung einer solchen Observablen liegt der Zustand im Eigenraum der zugeh"orenden
    Observablen. Da man die Richtung in diesem Eigenraum dadurch erh"alt, da\3 man den Ausgangszustand auf diesen Eigenraum projiziert, ergibt die Messung keine
    Information "uber die Richtung im Eigenraum. Je gr"o\3er die Entartung der gemessenen Observablen (entspricht der Dimension des Eigenraumes), um so weniger
    Inforamtion liefert die Messung "uber den Zustand des Systems. F"ur die Vorhersage bei so schwach pr"aparierten Zust"anden, m"ussen die bisherigen
    "Uberlegungen um statistische Methoden erg"anzt werden.
  \subsection{Der statisische Mischungszustand} 
    "Ahnlich wie in der klassischen Mechanik geht man in der Quantenmechanik vor, wenn man "uber das System nicht gen"ugend Informationen besitzt, um sagen zu 
    k"onnen, welche Richtung der Zustand des Quantensystems im Zustandsraum \raum{} einnimmt. 
    Sei $\Gamma$ ein Quantensystem. Die physikalischen Zust"ande des Systems $\Gamma$ entsprechen den Zust"anden im Zustandsraum \raum. Es soll jetzt das
    physikalische Verhalten des Systems $\Gamma$ mit statistischen Methoden beschrieben werden.
    \begin{itemize}
      \item[(i)]{\bf Statistischer Mischungszustand} \\
        Ein statistischer Mischungszustand $\Phi$ des Systems $\Gamma$ ist ein Tupel
        $$ \Phi = \left( \ket{\Phi_1}, p_1 ;\ket{\Phi_2}, p_2 ;\ldots \right) $$
        wobei $\ket{\Psi_i}$ Zust"ande in \raum{} sind, und $p_1 , p_2 ,\ldots \in \R$ mit $0 \le p_i \le 1$ f"ur alle $i$, und $\sum_i p_i = 1$.
        $p_m$ entspricht der Wahrscheinlichkeit, da\3 sich das System $\Gamma$ im Zustand $\ket{\Phi_m}$ befindet.
        Eine statistische Mischung $\Phi$ hei\3t ein reiner Zustand, wenn $p_{m_0} = 1$ f"ur ein festes $m_0$ und $p_m = 0$ f"ur alle $m \neq m_0$, sonst hei\3t
        $\Phi$ ein Mischungszustand.
      \item[(ii)]{\bf Messungen bei statistischen Mischungszust"anden} \\
        Es mu\3 gekl"art werden, was sich bei Vorliegen eines Mischungszustands $\Phi$ f"ur den Mittelwert $\mean{A}$ einer beliebigen Observablen $A$ ergibt. Nach den 
        statistischen Aussagen der Quantenmechanik ist der Mittelwert im Zustand $\ket{\Phi_m}$ gegeben durch
        $$ \mean{A}_m = \scalar{\Phi_m}{A}{\Phi_m} $$
        Diese Ergebnisse sind nun bei Vorliegen eines statistischen Mischungszustands mit den Gewichten $p_m$ zu gewichten, so da\3 man insgesamt f"ur den
        Mittelwert
        $\mean{A}_{\Phi}$ der Observablen im Mischungszustand $\Phi$ erh"alt
        $$ \mean{A}_{\Phi} = \sum_m p_m \mean{A}_m = \sum_m p_m \scalar{\Phi_m}{A}{\Phi_m}. $$
        Dabei ist zu beachten, da\3 man die Mittelung "uber die statistische Mischung in den Mittelwerten $\mean{A}_m$ ausf"uhrt. \\
        Entwickelt man die Zust"ande $\ket{\Phi_m}$ nach einer Basis $\menge{\ket{\Psi_i}}$, die zugleich die einfachen Einheitsvektoren des Operators A sind, so erh"alt man
        \begin{eqnarray*}
           \ket{\Phi_m} & = & \id\ket{\Phi_m} \\
                        & = & \sum_i \left(\ket{\Psi_i}\bra{\Psi_i}\right)\ket{\Phi_m} \\
                        & = & \sum_i \scalars{\Psi_i}{\Phi_m}\ket{\Psi_i}
        \end{eqnarray*}
        \begin{eqnarray*}
           \mean{A}_m  & = & \scalar{\Phi_m}{A}{\Phi_m} \\
                       & = & \sum_i \alpha_i \betrag{\scalars{\Phi_m}{\Psi_i}}^2 \\
           \mean{A}_{\Phi} & = & \sum_m p_m \left( \sum_i \alpha_i \betrag{\scalars{\Phi_m}{\Psi_i}}^2 \right)  
        \end{eqnarray*}
        In dieser Formel kommt ganz deutlich der Unterschied zwischen den beiden Arten der Mittelung zum Ausdruck: die quanten\-mechanische geschieht 
        mittels Wahrscheinlichkeitsamplituden und f"uhrt zu Interferenztermen, die Mittelung "uber die Mischung hingegen erfolgt durch die Gewichte $p_m$ und
        ergibt keine Interferenz.
    \end{itemize}


  \subsection{Das Konzept des Dichteoperators}
    Um die beiden verschiedenen Wahrscheinlichkeiten (klassische und quantenmechanische Wahrscheinlichkeiten) in einem gemeinsamen Kalk"ul zu verbinden, f"uhrt man
    den Dichteoperator
    $$ \rho = \sum_m p_m \ket{\Phi_m}\bra{\Phi_m} $$
    f"ur die statistische Mischung $\Phi$ ein. Bildet man in einer orthonormierten Basis $\menge{\ket{\Psi_i}}$ des Zustandraumes \raum{} die
    Matrixelemente des Dichteoperators $\rho$ 
    $$ \scalar{\Psi_k}{\rho}{\Psi_l} $$
    so erh"alt man die statistische Matrix oder die Dichtematrix (in Anlehnung an die Phasenraumdichte der statistischen Mechanik). Der Dichteoperator $\rho$ besitzt folgende 
    Eigenschaften (Verifikation durch Berechnung der Spur mit Hilfe einer beliebigen Basis): 
    \begin{eqnarray*}
      \mean{A}_{\Phi} & = & \Spur{\rho A} \\
      \Spur{\rho} & = & 1 \mbox{ da sich alle Wahrscheinlichkeiten $p_m$ zu 1 addieren }\\
      \Spur{\rho ^2} & = & 1 \Longleftrightarrow \mbox{ $\rho$ reiner Zustand} \\
      \Spur{\rho ^2} & \le & 1 \Longleftrightarrow \mbox{ $\rho$ statistischer Mischungszustand} \\
      \rho^{\dag} & = & \rho \mbox{ das hei\3t $\rho$ ist hermitesch}
    \end{eqnarray*}
    Besonders interessant ist die Formel $\mean{A}_{\Phi} = \Spur{\rho A}$, weil sie der Formel zu Berechnung des Mittelwertes einer Observablen in der klassischen 
    statistischen Physik entspricht, wobei die Integration "uber den Phasenraum der Spurbildung und die Wahrscheinlichkeitsdichte dem Dichteoperator entspricht.     
    Das Kriterium f"ur einen reinen oder gemischten Zustand ist $\Spur{\rho^2} = 1$ bzw. $\Spur{\rho^2} \le 1$
    F"ur einen Projektionsoperator $\Proj{\Psi}$ gilt
    $$ \Spur{\Proj{\Psi}\rho} = \sum_m p_m \betrag{\scalars{\Psi}{\Phi_m}}^2 $$
    ist also gleich der Wahrscheinlichkeit, den Zustand $\ket{\Psi}$ als Me\3ergebnis zu erhalten. Dies entspricht der Wahrscheinlichkeit, eine $1$ bei der Messung der Observablen
    $\Proj{\Psi}$ zu erhalten. Die Beschreibung der reinen Zust"ande mit Hilfe des Dichteoperators hat den Vorteil, da\3 der physikalisch unbedeutende Phasenfaktor verschwindet.
 
  \subsection{Besetzungen, Koh"arenzen}
    Was ist die physikalische Bedeutung der Matrixelemente $\rho_{kl}$ der Dichtematrix in der Basis $\menge{\ket{\Psi_i}}$, wobei die $\ket{\Psi_i}$ die Eigenvektoren einer
    vollst"andigen Observablen $A$ sind? \\
    F"ur die Hauptdiagonalelemente $\rho_{kk}$ gilt:
    $$ \rho_{kk} = \scalar{\Psi_k}{\rho}{\Psi_k} = \sum_m p_m \betrag{\scalars{\Psi_k}{\Phi_m}}^2 $$
    $\rho_{kk}$ entspricht also der Wahrscheinlickkeit, da\3 sich das Quantensystem nach der Messung von $A$ im Zustand ${\ket{\Psi_k}}$ befinden wird. Deswegen hei\3t $\rho_{kk}$
    die Besetzung des Zustands $\ket{\Psi_k}$. Wenn man diesselbe Messung $N$-mal unter denselben Bedingungen durchf"uhrt, wobei $N$ gro\3 ist, werden
    sich ann"ahrend $N\rho_{kk}$ Systeme im Zustand $\ket{\Psi_k}$ befinden. $\rho_{kk}$ ist eine positive reelle Zahl, die nur gleich Null ist, wenn alle
    $\betrag{\scalars{\Psi_k}{\Phi_m}}^2$ Null sind. \\
    F"ur das Matrixelement $\rho_{kl}$ ($k\neq l$) au\3erhalb der Hauptdiagonalen gilt:
    $$ \rho_{kl} = \sum_m p_m \scalars{\Psi_k}{\Phi_m}\overline{\scalars{\Psi_l}{\Phi_m}} $$
    Dieses Produkt beschreibt die Interferenzeffekte zwischen den Zust"anden $\ket{\Psi_k}$ und $\ket{\Psi_l}$, wenn der Zustand $\Phi_m$ eine
    koh"arente lineare Superposition dieser Zust"ande ist. $\rho_{kl}$ ist also der Mittelwert  der Interferenzterme aller m"oglichen 
    Zust"ande des statistischen Mischungszustandes. Im Gegensatz zu den Besetzungen kann $\rho_{kl}$ Null sein, obwohl keines der Produkte Null ist.
    $\rho_{kk}$ ist eine Summe von nicht negativen reellen Zahlen, w"ahrend $\rho_{kl}$ eine Summe von komplexen Zahlen ist. Wenn $\rho_{kl}$ Null ist, dann
    l"oschen sich durch die statistische Mittelwertbildung alle Interferenzeffekte zwischen $\ket{\Psi_k}$ und $\ket{\Psi_l}$ aus. Wenn dagegen $\rho_{kl}$
    nicht Null ist, dann existiert noch eine gewisse Koh"arenz zwischen diesen Zust"anden. Deswegen werden diese Matrixelemente $\rho_{kl}$ Koh"arenzen 
    genannt. \\
    Die Unterscheidung zwischen Besetzungen und Koh"arenzen h"angt offensichtlich von der Basis $\menge{\Psi_i}$ ab. Da $\rho$ hermitesch ist,
    kann man eine orthonormale Basis $\menge{\ket{\Pi_i}}$ finden in der $\rho$ diagonal ist.
    $$ \rho = \sum_m \pi_m \Proj{\Pi_m} $$
    Da $\rho$ positiv semidefinit ist, und $\Spur{\rho} = 1$ gilt:
    $$ 0 \le \pi_m \le 1 $$
    $$ \sum_m \pi_m = 1 $$
    $\rho$ beschreibt also eine statistische Mischung zwischen den Zustanden $\ket{\Pi_m}$ mit den Wahrscheinlichkeiten $\pi_m$. Es bestehen keine
    Koh"arenzen zwischen diesen Zust"anden, da $\rho$ in dieser Basis diagonal ist. 
  \subsection{Beispiel: Polarisierte Photonen}
    Betrachten wir als Beispiel polarisierte Photonen, die auf einen Analysator treffen, um die Unterschiede zwischen reinen und statischischen Mischungszust"ande und die Bedeutung 
    der linearen Superpositon zu verdeutlichen. Die lineare Superposition 
    $$ \ket{\Psi} = \frac{1}{\sqrt{2}}\ket{0} + \frac{1}{\sqrt{2}}\ket{1} $$
    entspricht einem linear polarisierten Photon. Bei dem statistischen Mischungszustand 
    $$ \Phi = \left( \ket{0},\frac{1}{2};\ket{1},\frac{1}{2}\right) $$
    ist die H"alfte der emittierten Photonen $\ket{0}$-polarisiert und die andere H"alfte $\ket{1}$-polarisiert. Um das Verhalten des Systems bei Vorliegen von $\ket{\Psi}$ bzw.
    von $\Phi$ zu studieren, berechnet man die Wahrscheinlichkeiten, da\3 das Photon durch einen im Strahlengang positionierten Analysator durchgelassen wird. \\
    \begin{tabular}{ccrcl}
      Messung von      &    in Richtung     & \multicolumn{3}{l}{mit Wahrscheinlichkeit}\\
      $\ket{\Psi}$     &   $\ket{0}$        &   $\betrag{\scalars{0}{\Psi}}^2$ &=& $\frac{1}{2}$ \\
      $\ket{\Psi}$     &   $\ket{1}$        &   $\betrag{\scalars{1}{\Psi}}^2$ &=& $\frac{1}{2}$ \\
      $\Phi$           &   $\ket{0}$        &   $\frac{1}{2}\betrag{\scalars{0}{0}}^2+\frac{1}{2}\betrag{\scalars{0}{1}}^2$ &=& $\frac{1}{2}$ \\
      $\Phi$           &   $\ket{1}$        &   $\frac{1}{2}\betrag{\scalars{1}{0}}^2+\frac{1}{2}\betrag{\scalars{1}{1}}^2$ &=& $\frac{1}{2}$ \\
    \end{tabular} \\
    Die Wahrscheinlichkeiten sind bei den jeweiligen Messungen gleich. Es stellt sich also die Frage, ob $\ket{\Psi}$ als $\Phi$ interpretiert werden kann. Man sieht, da\3 
    $\ket{\Psi}$ und $\Phi$ grunds"atzlich unterschiedliche physikalische Sitiuationen beschreiben, wenn man eine Messung in Richtung $\ket{\Psi}$ durchf"uhrt: \\
    \begin{tabular}{cl}
       Messung von    &    Wahrscheinlichkeit \\
       $\ket{\Psi}$   &    $\betrag{\scalars{\Psi}{\Psi}}^2=1$ \\
       $\Phi$         &    $\frac{1}{2}\betrag{\scalars{\Psi}{0}}^2 + \frac{1}{2}\betrag{\scalars{\Psi}{0}}^2 = \frac{1}{2}\frac{1}{2} + \frac{1}{2}\frac{1}{2} = \frac{1}{2}$
    \end{tabular} \\
    Man kann diese Informationen viel einfacher aus den zugeh"origen Dichtematrizen ablesen.\\
    Der Dichteoperator f"ur die lineare Superposition $\ket{\Psi}$ in der Basis $\menge{\ket{0},\ket{1}}$ ist
    $$ \rho_{\ket{\Psi}} = \Proj{\Psi} = \frac{1}{2}\Proj{0} + \frac{1}{2}\Proj{1} + \frac{1}{2}\ket{0}\bra{1} + \frac{1}{2}\ket{1}\bra{0}. $$
    Die ersten beiden Summanden sind die Besetzungen, w"ahrend die letzten beiden die Koh"arenzen beschreiben. Die Dichtematrix von $\rho_{\ket{\Psi}}$ ist
    $$ \def\arraystretch{1.5}       \rho_{\ket{\Psi}} =  \left(
          \begin{array}{cc}                                                    
            \frac{1}{2}&\frac{1}{2}\\
            \frac{1}{2}&\frac{1}{2}
          \end{array}
        \right) .
    $$
    Der Dichteoperator f"ur den statistischen Mischungszustand $\Phi$ in der Basis $\menge{\ket{0},\ket{1}}$ ist
    $$ \rho_{\Phi} = \frac{1}{2}\Proj{0} + \frac{1}{2}\Proj{1}. $$
    Die Dichtematrix von $\rho_{\Phi}$ ist
    $$ \rho_{\Phi} = \left(  
         \begin{array}{cc} 
            \frac{1}{2}&0\\
            0&\frac{1}{2}
         \end{array}
       \right) .
    $$
    Jetzt wird eine Basis gesucht, in der der Dichteoperator $\rho_{\ket{\Psi}}$ diagonal ist, da\3 es also keine Koh"arenzen zwischen den Basiszust"anden gibt.
    Die Eigenwerte von $\rho_{\ket{\Psi}}$ sind $0$ und $1$, da dieser Operator ein Projektionsoperator ist. Die zugeh"origen Eigenvektoren sind
    $\frac{1}{\sqrt{2}}\left( \ket{0}+\ket{1}\right)$ und $\frac{1}{\sqrt{2}}\left( \ket{0}+\ket{1}\right)$. Der Basiswechsel wird durch die Hadamardmatrix
    $$ \frac{1}{\sqrt{2}} \left(\matrix{1&1\cr
                                 1&-1}\right) $$
    realisiert. Die Dichtematrix von $\rho_{\ket{\Psi}}$ in der neuen Basis ist
    $$ \left(\matrix{1&0\cr
                     0&0}\right). $$
    Die Dichtematrix von $\rho_{\Phi}$ in der neuen Basis ist
    $$ \left(\matrix{\frac{1}{2}&0\cr
                     0&\frac{1}{2}}\right). $$
    Einer Messung in anderen Richtung entpricht die Betrachtung des Dichteoperators in einer anderen Basis. \\
    Eine weitere interessante Beobachtung ist die, da\3 die statistischen Mischungszust"ande werden eindeutig durch die Dichteoperatoren charakterisiert werden, denn
    $$ \Phi_1 = \left( \ket{0}, \frac{1}{2}; \ket{1}, \frac{1}{2}\right) $$
    $$ \Phi_2 = \left( \ket{0}, \frac{1}{3}; \frac{1}{2}\ket{0}+\frac{\sqrt{3}}{2}\ket{1},\frac{1}{3}; \frac{1}{2}\ket{0}-\frac{\sqrt{3}}{2}\ket{1},\frac{1}{3} \right) $$
    haben denselben Dichteoperator
    $$ \rho = \frac{1}{2}\ket{0} + \frac{1}{2}\ket{1}. $$
    Diese beiden statistischen Mischungszust"ande k"onnen mit keiner Messung unterschieden werden.
   \subsection{Dynamik eines statistischen Mischungszustandes}
    F"ur die Behandlung der Dynamik eines quantenmechanischen Systems $\Gamma$, das durch einen statistischen Mischungszustand charakterisiert ist, ist es 
    notwendig, auch f"ur den Dichteoperator $\rho$ ein Bewegungsgesetz anzugeben, das dann erlaubt, die Zeitabh"angigkeit von Erwartungswerten ebenso zu 
    berechnen wie im reinen Fall. Der Unterschied gegen"uber dem reinen Zustand kommt in den Gewichten $p_m$ zum Ausdruck, die die fehlende Information "uber
    den im Gemisch tats"achlich vorliegenden Zustand  beschreiben. Weil in der quantenmechanischen Dynamik das System sich selbst "uberlassen ist und keine 
    Messungen an dem System durchgef"uhrt werden, "andert sich die Information "uber den Zustand nicht. Die $p_m$ sind daher keine Zeitfunktionen.
    Die zeitliche Entwicklung der m"oglichen Zustande $\menge{\ket{\Phi_m}}$ des statistischen Mischungszustandes $\Phi$ wird durch die Schr"odingergleichung  
    $$ i\hbar \difft \ket{\Phi_m\left( t\right)} = H\ket{\Phi_m\left( t\right)} $$
    beschrieben. Dabei sind $\menge{\ket{\Phi_m}}$ die m"oglichen Zust"ande zum Zeitpunkt $t=0$. Wenn man die Schr"odingergleichung konjugiert, erh"alt man die 
    entsprechende Gleichung f"ur $\menge{\bra{\Phi_m}}$ 
    $$ -i\hbar \difft \bra{\Phi_m\left( t\right)} = \bra{\Phi_m\left( t\right)}H^{\dag}. $$
    Mit Hilfe der Produktregel erh"alt man die zeitliche Entwicklung der Projektionsoperatoren $\Proj{\Phi_m}$
    \begin{eqnarray*}
      i\hbar \difft\left( \Proj{\Phi_m} \right)  & = & \left( i\hbar \difft\ket{\Phi_m} \right)\bra{\Phi_m} + \ket{\Phi_m} \left( i\hbar \difft\bra{\Phi_m}\right) \\
                                                 & = & i\hbar\left( H\Proj{\Phi_m} - \Proj{\Phi_m}H^{\dag}\right) \\
                                                 & = & i\hbar\left( H\Proj{\Phi_m} - \Proj{\Phi_m}H\right) \\
                                                 & = & i\hbar\kommutator{H}{\Proj{\Phi_m}}.
    \end{eqnarray*}
    F"ur die zeitliche Entwicklung des Dichteoperators $\rho$ ergibt sich damit
    $$ i\hbar\difft \rho\left( t\right) = \kommutator{H\left( t\right)}{\rho\left( t\right)}. $$
    Diese Gleichung hei\3t die von Neumann-Gleichung; sie ist das quantenmechanische Analogon der Liouville-Gleichung der klassischen statistischen Mechanik, die die zeitliche 
    Entwicklung der Wahrrscheinlichkeitsdichte im Phasenraum beschreibt. \\
    $\Spur{\rho^2}$ ist zeitunabh"angig. Also bleibt ein reiner (gemischter) Zustand rein (gemischt).
    \begin{eqnarray*}
      \Spur{\rho^2\left( t\right)} & = & \Spur{U\rho\left( t_0 \right) U^{\dag} U\rho\left( t_0 \right) U^{\dag}} \\
                                   & = & \Spur{U\rho^2\left( t_0\right) U^{\dag}} \\
                                   & = & \Spur{\rho^2 \left( t_0\right)} 
    \end{eqnarray*}
    wegen der Invarianz unter Konjugation der Spur.

    \begin{thebibliography}{xxxxxxxxxx}
    \bibitem[Cohen]{Cohen}
      Cohen-Tannoudji, {\em Quantum Mechanics}, Wiley, 1977
    \bibitem[Fick]{Fick}
      Eugen Fick, {\em Einf"uhrung in die Grundlagen der Quantentheorie}, Aula-Verlag, 1984
    \bibitem[Sakurai]{Sakurai}
      J.J. Sakurai, {\em Modern Quantum Mechanics}, Addison-Wesley, 1985
    \bibitem[Schwabl]{Schwabl}
      Franz Schwabl, {\em Quantenmechanik}, Springer, 1992
    \bibitem[Zeidler]{Zeidler}
      Eberhard Zeidler, {\em Applied Functional Analysis}, Springer, 1995
    \end{thebibliography}

%% \end{document}      



 

%%
%% Bemerkungen zur Ausarbeitung
%%
%% - erster Befehl:
%%
%%   \Vortrag{Titel}{Autor}
%%
%% - oberste Gliederungsstufe ist \section
%%
%% - \begin{document} und \end{document} entfallen
%%
%% - eigenes Literaturverzeichnis
%%

\end{document}


